\scanpage{156}
\mspage{72}
% archive sheet number: 155

pour tout $x$ dans $M$ on a

\[
x=\underset{s\in S_{/x}}{\underrightarrow{\lim}}\ s
\tag*{(4.3.6.2)}
\]

cf SGA 4 I \unsure{7|?}.2. \unsure{(i)|bis}). Soit $\pi_0=\mathrm{card}\,\mathrm{Fl}(S)$, \deleted{et soit} et pour tt cardinal $\pi\geq\pi_0$, soit

\[
\begin{aligned}
\mathrm{Filt}^{\pi}M^{\deleted{\unsure{\pi}}}
&=\text{sous-catégorie pleine de }M\text{ formée}\\
&\phantom{{}={}}\text{des objets }x\text{ tels que }\mathrm{card}(\mathrm{Fl}\,S_{/x})\leq\pi.
\end{aligned}
\tag*{(4.3.6.3)}
\]

Je veux prouver que \deleted{\unsure{ceci définit}} \inserted{l'on \unsure{trouve|s'trouve} ainsi} une filtration cardinale sur $M$, i.e. que les conditions a) b) c) de SGA 4 I 9.12. sont satisfaites.

C'est prouvé dans loc. cit. 9.13. sous des \deleted{conditions} hypothèses légèrement différentes, la condition d'accessibilité de $M$ y est remplacée par celle \deleted{que} (beaucoup moins sympa !) que les familles de flèches épimorphiques strictes dans $M$ sont épimorphiques strictes universelles. Cette condition n'est utilisée que dans la démonstration de \inserted{la} condition \unsure{b)}, à savoir la stabilité de $\mathrm{Filt}^\pi M$ par \deleted{\unsure{petites}} \inserted{$\underset{I}{\underrightarrow{\lim}}$} de cardinal $\leq\pi$ (sans avoir d'ailleurs à supposer $I$ filtrante, ni a fortiori grand devant $\pi_0$).

Il faut d'ailleurs dire que la définition de $\mathrm{Filt}^\pi M$ donnée dans loc cit. n'est pas tout à fait la $\unsure{=|\simeq}$ que (4.3.6.3),

